\section{本讲小结}

综合成像按实际经过的界面逐段处理，每次重新判断发散、会聚和量距方向。平行板在近轴成像传递中贡献 $t/n$，法向光程却为 $nt$；往返系统的镜间传播累计两次，透镜作用也保留两次。银层所在的曲面改变实际界面顺序，因此等效光焦度不能只由透镜名称判断。

清晰观察与视角放大是不同的条件。先求最终像的坐标，再从眼位量观察距离，并按实际近点和远点判断调焦。视角放大率以同一物体在约定比较距离处的裸眼视角为基准；若实际近点不同，或眼睛与目镜有间隔，应把这些距离分别写入几何关系。

望远镜的无焦条件连接物镜后焦平面与目镜物方焦平面，角放大率的符号由实际出射倾角决定。开普勒式在两镜间形成真实中间像，伽利略式以负目镜接收尚未完成会聚的虚物。显微镜先由物镜作横向放大，再由目镜作视角放大；固定镜距下改变最终观察状态，还会改变物镜的物距与放大率。

出瞳是孔径光阑经后续光组形成的像。其位置、直径与眼瞳共同决定通光，不能用远物的横向成像代替孔径成像，也不能只由角放大率断定视场全部可见。设计空心光束时还应区分球面反射处的精确几何与末端透镜的近轴条件，逐项检查传播支路、会聚目标、环半径与实际孔径。
